\chapter{Business Operations}

\section{Architecture: Two Servers}

Connector separates execution from control. The \textbf{Platform Server} runs
on the customer's infrastructure and executes agents; the
\textbf{License/Control Server} runs on Connector's hosted infrastructure and
manages accounts, licensing, and billing. The two servers communicate only
for authentication and telemetry — no agent data or MemPackets ever leave the
customer's node.

\begin{figure}[H]
\centering
\begin{tikzpicture}[
  server/.style={rectangle, rounded corners=6pt, draw=#1!70, fill=#1!8,
                 text width=5.4cm, align=center, minimum height=4.2cm,
                 font=\small, inner sep=8pt},
  feat/.style={rectangle, rounded corners=2pt, draw=#1!40, fill=#1!5,
               text width=4.6cm, align=center, minimum height=0.65cm,
               font=\footnotesize, inner sep=4pt},
  node distance=0.25cm
]

%% Platform Server
\node[server=CKernel] (plat) at (-3.4,0) {};
\node[font=\bfseries\small, text=CKernel] at (-3.4, 1.8) {Platform Server};
\node[font=\footnotesize\itshape, text=gray] at (-3.4, 1.42) {Customer infrastructure};

\node[feat=CKernel] (f1) at (-3.4, 0.9)  {Agent execution \& memory kernel};
\node[feat=CKernel] (f2) at (-3.4, 0.18) {connectorctl CLI + REST API};
\node[feat=CKernel] (f3) at (-3.4,-0.54) {Local dashboard (API-key auth)};
\node[feat=CKernel] (f4) at (-3.4,-1.26) {Prometheus + structured logs};

%% License/Control Server
\node[server=CBiz] (ctrl) at (3.4,0) {};
\node[font=\bfseries\small, text=CBiz] at (3.4, 1.8) {License/Control Server};
\node[font=\footnotesize\itshape, text=gray] at (3.4, 1.42) {Connector-hosted SaaS};

\node[feat=CBiz] (g1) at (3.4, 0.9)  {Customer portal (signup, login)};
\node[feat=CBiz] (g2) at (3.4, 0.18) {API key issuance \& management};
\node[feat=CBiz] (g3) at (3.4,-0.54) {License activation \& enforcement};
\node[feat=CBiz] (g4) at (3.4,-1.26) {Stripe billing \& usage telemetry};

%% Arrows
\draw[darr=gray!50]
  (plat.east) -- node[above, font=\footnotesize\itshape]{license check}
               node[below, font=\footnotesize\itshape]{usage report}
  (ctrl.west);

\node[below=0.6cm of plat.south, font=\footnotesize\itshape, text=gray]
  {Agent data never leaves this server};

\end{tikzpicture}
\caption{Two-server architecture — execution is local; control and billing are hosted}
\label{fig:two-server}
\end{figure}

This separation is a deliberate \textbf{data sovereignty design}:
compliance-sensitive organisations (healthcare, finance) can run agent workloads
on-premise while still benefiting from the hosted portal, admin panel, and
billing infrastructure.

\section{Commercial Thesis}

Connector should be sold as a \textbf{governed runtime and evidence layer for
production AI}, not as a replacement for every orchestration, storage,
observability, and compliance tool in the stack. The commercial message must
stay narrow at entry and only broaden after trust is earned.

\begin{warnbox}{Avoid the platform-first trap}
Do not lead with ``Connector replaces your whole AI stack.'' The better entry
message is:

\smallskip
\centering
\textbf{``When an AI workflow goes wrong, Connector shows what happened, why it
happened, which policy allowed it, and which receipt proves it.''}
\end{warnbox}

\subsection{Why now}

The 2025--2026 market is favourable for a focused infrastructure vendor that
helps enterprises move AI from experimentation into governed production.
Menlo Ventures reports that enterprises now \textbf{buy more AI than they
build}, with purchased solutions materially outpacing internal builds, and that
AI deals reach production at higher rates than traditional SaaS. At the same
time, investor guidance for 2026 remains traction-first: paid pilots, capital
efficiency, and a credible next-round plan matter more than ambitious platform
narratives.

\section{Product Positioning}

\begin{conceptbox}{Core product definition --- do not change this}
\centering
\large\bfseries Connector = AI Decision Audit, Proof, and Cost Control Layer

\smallskip
\normalsize\normalfont\itshape
``When your AI makes a decision, Connector shows what happened, why it happened,\\
proves it cryptographically, and controls the cost --- in real time.''
\end{conceptbox}

\subsection{Three Pillars}

Three capabilities. Not a platform narrative. Not a suite. Three specific
things a buyer can hold in their head and explain to their own team.

\begin{table}[H]
\centering
\renewcommand{\arraystretch}{1.7}
\begin{tabularx}{\textwidth}{@{}clXX@{}}
\toprule
\textbf{\#} & \textbf{Pillar} & \textbf{What it means} & \textbf{Why a buyer pays for it} \\
\midrule
1 & \textbf{Audit} & Every agent decision produces a trace: what
  happened, in what order, with what data, at what time. Tamper-evident.
  Kernel-level, not app-level. & ``When something went wrong, I can replay
  exactly what the agent did and why. I do not need a developer in the room
  to explain it.'' \\
\midrule
2 & \textbf{Proof} & Each trace is signed with a cryptographic receipt
  (CID-addressed, verifiable). The signature cannot be faked or retroactively
  edited. Structured for compliance output: HIPAA, SOC2~CC6/CC7, FCA. & ``I
  can show this to my auditor, my regulator, or my legal team. It is not a
  log. It is signed evidence.'' \\
\midrule
3 & \textbf{Cost Control} & Every agent action records token usage,
  compute spend, and API calls. Hard caps are enforced at the kernel before
  the action executes, not after. Budget per agent, per session, per team. &
  ``I know exactly what each AI workflow costs. I can set a hard limit and
  it will not be exceeded. No surprise bills.'' \\
\bottomrule
\end{tabularx}
\caption{The three pillars of Connector --- Audit, Proof, Cost Control}
\end{table}

\subsection{Entry Wedge}

The correct wedge is not ``replace all AI infrastructure.'' It is to land on
one specific painful workflow, deliver all three pillars against it, and
then expand.

\begin{figure}[H]
\centering
\begin{tikzpicture}[
  box/.style={rectangle, rounded corners=4pt, draw=#1!70, fill=#1!8,
              text width=2.9cm, align=center, minimum height=1.0cm,
              font=\footnotesize\bfseries, inner sep=5pt},
  node distance=1.0cm
]
\node[box=CDanger] (w1) {Trace one\\broken decision};
\node[box=CLogic, right=of w1] (w2) {Prove why\\it happened};
\node[box=CExec, right=of w2] (w3) {Control cost\\and budget};
\node[box=CBiz, right=of w3] (w4) {Expand to\\all workflows};
\node[box=CKernel, right=of w4] (w5) {Runtime of\\record};

\foreach \a/\b in {w1/w2,w2/w3,w3/w4,w4/w5}
  \draw[arr=gray!60] (\a) -- (\b);

\end{tikzpicture}
\caption{Commercial wedge --- land on one broken decision, expand to runtime of record}
\label{fig:commercial-wedge}
\end{figure}

\section{Who Actually Pays: Reality-First ICP}

\begin{conceptbox}{First principle --- lock this}
\centering
\textbf{Not} ``AI startups.'' \textbf{Not} ``customer support SaaS.''

\smallskip
\textbf{Teams where AI is already making decisions that affect money or customers
--- and those decisions are going wrong or unclear right now.}

\smallskip
\textit{If they are not feeling it weekly, they will not open their wallet.}
\end{conceptbox}

\subsection{The Three-Question Qualifier (Use Before Every Call)}

All three must be true or move on:

\begin{enumerate}[leftmargin=1.8em, itemsep=3pt]
  \item \textbf{AI is live in production} with real users making real decisions ---
    not a test, not a prototype, not ``we are exploring.''
  \item \textbf{Decisions affect money or customer trust} --- refunds issued,
    transactions blocked, tickets actioned, approvals made automatically.
  \item \textbf{They have already felt it go wrong} --- a wrong refund they
    could not explain, a cost spike, an enterprise client who asked for an
    audit log they could not produce.
\end{enumerate}

\begin{warnbox}{Disqualify fast --- your time is the constraint in 2026}
No prior AI incident = no urgent budget. Ask in the first 2 minutes:
\textit{``Has your AI ever made a decision you could not explain to a customer
or your team?''} No = end the call. Yes = keep going.
\end{warnbox}

%%───────────────────────
\subsection{Segment 1: Ecommerce AI Support \texorpdfstring{---}{-} Enter Here First}
%%───────────────────────

\textbf{Exact type:} Shopify-focused SaaS, DTC automation tools, helpdesk
products with AI automation. Companies whose product \emph{is} AI-driven
customer ops --- refunds, returns, order edits, discount decisions.

\begin{table}[H]
\centering\renewcommand{\arraystretch}{1.55}
\begin{tabularx}{\textwidth}{@{}lX@{}}
\toprule
\textbf{Dimension} & \textbf{Concrete detail} \\
\midrule
Real pain & ``We lost \$4K this week from wrong refunds.'' ``AI is too
  aggressive on discounts.'' ``A customer escalated and we have no idea
  why it issued that refund.'' Happens weekly. \\
Who pays & Founder or Head of Ops (feels the financial pain). CX lead
  (faces escalations). CTO signs. \\
Pitch & \textit{``Your AI is issuing refunds you cannot explain. We give
  you the full trace, the reason, and the proof of every decision --- so
  you stop losing money and explain instantly.''} \\
Demo & AI issued a \$120 refund incorrectly. Show: trace (what happened),
  explain (which rule applied), prove (signed receipt), cost (token usage).
  5 minutes. One bad decision. Buyer understands immediately. \\
Sales cycle & 2--4 weeks. One buyer. No compliance review. High urgency. \\
\midrule
\textbf{Why \$10K--\$20K / 6 months is easy to approve} &
  Processing 500 refunds/month at \$80 average = \$40K/month in refund
  volume. Even 3\% wrong = \$1,200/month of bad money out. Debugging each
  incident manually = 3--4 engineer hours at \$150/hr $\times$ 15 incidents =
  \$6,750/month hidden cost. Total problem cost: \$8K--\$20K/month. Connector
  at \$2K--\$3K/month to fix it. The maths approves itself in 10 minutes. \\
\bottomrule
\end{tabularx}
\caption{Segment 1 --- Ecommerce AI Support: best wedge, fastest close, clearest ROI}
\end{table}

%%───────────────────────
\subsection{Segment 2: AI Support Agent Startups \texorpdfstring{---}{-} Action-Takers Not Chatbots}
%%───────────────────────

\textbf{Exact type:} Startups building agents that \emph{take actions} ---
close tickets, call APIs, issue credits, resolve workflows end-to-end. Not
chatbots that answer questions. Agents that do things with consequences.

\begin{table}[H]
\centering\renewcommand{\arraystretch}{1.55}
\begin{tabularx}{\textwidth}{@{}lX@{}}
\toprule
\textbf{Dimension} & \textbf{Concrete detail} \\
\midrule
Real pain & ``Customer complained the AI did the wrong thing and we spent
  two days figuring out why.'' ``Our biggest enterprise prospect asked for
  an audit log and we have nothing to show.'' ``We keep getting blocked in
  procurement over AI compliance questions.'' \\
Who pays & CTO or Founder. Head of Product when an enterprise deal is at risk. \\
Pitch & \textit{``Your enterprise customers will ask `why did your AI do this?'
  before they sign. We give you that answer instantly, with a signed receipt
  you can put in front of their compliance team.''} \\
Sales cycle & 3--6 weeks. CTO approves directly. One blocked enterprise deal
  creates immediate urgency. \\
\midrule
\textbf{Why \$10K--\$20K / 6 months is easy to approve} &
  One blocked or lost enterprise deal due to lack of auditability = \$30K--\$150K
  ARR delayed or gone. Connector at \$3K/month unblocks that deal in the next
  sales cycle. 1 deal recovered = 5--10$\times$ the pilot cost. Additionally:
  the audit trace becomes a \emph{sales feature} they show in their demo.
  ``Here is the signed proof of every decision our agent made.'' That line
  closes deals. \\
\bottomrule
\end{tabularx}
\caption{Segment 2 --- AI Support Agent startups: audit as sales feature and deal-unlocker}
\end{table}

%%───────────────────────
\subsection{Segment 3: Fintech \texorpdfstring{---}{-} Selective, Not All}
%%───────────────────────

\textbf{Exact type:} Fraud detection tools, payment monitoring SaaS, risk
scoring platforms. \textbf{Not banks. Not early-stage. Not ``exploring
AI.''} Only startups where AI is already making real financial decisions
in production today.

\begin{table}[H]
\centering\renewcommand{\arraystretch}{1.55}
\begin{tabularx}{\textwidth}{@{}lX@{}}
\toprule
\textbf{Dimension} & \textbf{Concrete detail} \\
\midrule
Real pain & ``AI blocked a legitimate transaction and the customer left.''
  ``We need explainability for our next compliance review.'' ``Regulators
  are asking how our AI makes decisions and we have no structured answer.'' \\
Who pays & CTO (primary). Head of Risk. \\
Pitch & \textit{``If your AI blocks or approves money, you need to prove
  why --- to regulators, to customers, and to your own risk team. We give
  you a signed, verifiable record of every decision.''} \\
Sales cycle & 4--10 weeks. Slower than Segment~1 but CTO can approve
  under \$10K directly. Compliance officer adds 2--3 weeks. \\
\midrule
\textbf{Why \$10K--\$20K / 6 months is easy to approve} &
  One wrong AI decision resulting in a regulatory inquiry = legal review
  at \$300--\$600/hr, 10--40 hours minimum = \$3K--\$24K per incident.
  Manual compliance analyst reviewing AI decisions = 0.5 FTE at \$80K/year
  = \$3,300/month. Connector at \$3K/month replaces or eliminates both. \\
\bottomrule
\end{tabularx}
\caption{Segment 3 --- Fintech: regulatory urgency creates real forcing function; enter after 2 support wins}
\end{table}

%%───────────────────────
\subsection{Segment 4: Internal AI Ops \texorpdfstring{---}{-} Only After Incidents}
%%───────────────────────

\textbf{Exact type:} Companies using AI internally for approvals, workflows,
or ops automation. \textbf{Only approach if they have already had an incident}
--- an AI action that cost them time, money, or internal credibility.
Without an incident, urgency is low and the deal will stall.

\begin{table}[H]
\centering\renewcommand{\arraystretch}{1.55}
\begin{tabularx}{\textwidth}{@{}lX@{}}
\toprule
\textbf{Dimension} & \textbf{Concrete detail} \\
\midrule
Real pain & ``AI did something unexpected last month and it took us three
  days to untangle.'' ``Cost is completely unpredictable.'' ``We don't fully
  trust it yet but we're running it anyway.'' \\
Close probability & Low without a prior incident. If they can name a specific
  incident that cost them, it becomes medium. \\
Strategy & Do not outbound to this segment. If they come inbound or if they
  mention a specific incident, qualify and proceed. Otherwise skip. \\
\bottomrule
\end{tabularx}
\caption{Segment 4 --- Internal AI Ops: only viable if they have had a named incident}
\end{table}

%%───────────────────────
\subsection{Who Will Not Pay --- Save Your Time}
%%───────────────────────

\begin{warnbox}{Do not contact these company types --- they will never close in 2026}
\begin{itemize}[noitemsep, leftmargin=1.4em]
  \item \textbf{Chatbot SaaS.} They answer questions, they do not take actions.
    No financial consequence from a bad response = no urgency to audit it.
  \item \textbf{Early AI startups with no production usage.}
    They will say ``this is interesting, come back in six months.''
    They are not lying. They mean it. Move on.
  \item \textbf{Enterprises without active AI decisions.}
    Large companies exploring AI take 6--18 months from pilot approval to
    contract. You do not have that runway in 2026.
  \item \textbf{``Exploring AI'' companies.}
    If their website says ``we are integrating AI'' or ``AI-powered coming
    soon'' --- they are 3--6 months from feeling the pain. You want companies
    that already feel it.
\end{itemize}
\end{warnbox}

%%───────────────────────
\subsection{The Pilot Offer --- Say This Exactly}
%%───────────────────────

\begin{conceptbox}{The offer --- one workflow, 30 days, sharp scope}
\centering
\textit{``Give us one AI workflow where decisions matter --- refunds,
approvals, or actions. In 30 days we will give you:}

\smallskip
\begin{itemize}[noitemsep, leftmargin=2em]
  \item \textit{Full traceability of every decision in that workflow}
  \item \textit{Explanation of every action the agent took and why}
  \item \textit{Proof --- a signed receipt you can show your compliance team}
  \item \textit{Cost visibility and hard budget control per session}
\end{itemize}

\smallskip
\textit{One workflow. \$2K--\$5K/month. Six months.
You will know in 30 days whether it is worth continuing.''}
\end{conceptbox}

The 30-day promise is critical. It converts a 6-month commitment into a
30-day decision. The buyer does not need to commit to the full value ---
they commit to seeing one workflow traced. If it works (it will), they
renew. You get 6 months of revenue from a 30-day promise.

%%───────────────────────
\subsection{How to Find the Right Companies}
%%───────────────────────

\begin{table}[H]
\centering\renewcommand{\arraystretch}{1.55}
\begin{tabularx}{\textwidth}{@{}lXl@{}}
\toprule
\textbf{Signal} & \textbf{What it means} & \textbf{Strength} \\
\midrule
Website says ``AI auto-resolves'' or ``AI handles X'' &
  AI is in production taking actions & Strong \\
Job posting for LLM Engineer or AI Engineer &
  They are actively building or scaling AI in production & Strong \\
Case study: ``We reduced support cost by X\% using AI'' &
  AI is live, decisions are happening, mistakes exist somewhere &
  \textbf{Best signal} \\
Product demo shows AI taking autonomous actions &
  They have an action-taking agent, not just a chatbot & Strong \\
LinkedIn: employees with title ``AI Ops'' or ``Agent Engineer'' &
  Company is past exploration phase & Medium \\
Trustpilot or G2 reviews mention ``AI did wrong thing'' &
  Public evidence of live AI decisions going wrong & Very strong \\
\bottomrule
\end{tabularx}
\caption{Prospecting signals --- how to identify companies where AI is already causing problems}
\end{table}

\begin{insightbox}{The 20-5-1 outreach formula for the first pilot}
Find 20 companies matching Segment~1 signals. Shortlist 5 that show the
``best signal'' (reduced cost by X using AI, or a public review about wrong
AI action). Pitch only this one question: \textit{``Are you able to explain
every AI refund decision today?''} If the answer is no or hesitation ---
schedule the demo. Target: 1 pilot closed from 20 companies contacted.
That is the realistic close rate. If you contact fewer than 20, you are
not doing outreach --- you are waiting.
\end{insightbox}

\begin{conceptbox}{The only metric that matters in outreach}
\centering
\textbf{``Don't sell Connector. Sell stopping AI mistakes where money is leaking.''}

\smallskip
\textit{You need 3 companies where AI is already hurting them. Not 100 customers.
Not a big market. Three companies. Find them. Close them. That is 2026.}
\end{conceptbox}

\section{Customer Signup and Login}

\subsection{Signup Flow}

\begin{figure}[H]
\centering
\begin{tikzpicture}[
  stepnode/.style={rectangle, rounded corners=3pt, draw=CBiz!60, fill=CBiz!7,
               text width=3cm, align=center, minimum height=0.85cm,
               font=\footnotesize\bfseries, inner sep=5pt},
  dec/.style={diamond, draw=CLogic!60, fill=CLogic!7,
              text width=1.8cm, align=center, minimum height=0.9cm,
              font=\scriptsize, inner sep=3pt},
  node distance=0.5cm
]
\node[stepnode] (visit)  {Visit portal};
\node[stepnode, right=of visit]  (form)  {Submit\\registration};
\node[stepnode, right=of form]   (email) {Verify\\email};
\node[stepnode, right=of email]  (dash)  {Portal\\dashboard};
\node[stepnode, right=of dash]   (key)   {Create\\API key};
\node[stepnode, right=of key]    (node)  {Activate\\license};

\foreach \a/\b in {visit/form,form/email,email/dash,dash/key,key/node}
  \draw[arr=CBiz!60] (\a) -- (\b);

\node[below=0.3cm of form, font=\scriptsize\itshape, text=gray, align=center]
  {Argon2 password\\hash};
\node[below=0.3cm of email, font=\scriptsize\itshape, text=gray, align=center]
  {HMAC-signed\\token};
\node[below=0.3cm of key, font=\scriptsize\itshape, text=gray]
  {shown once};
\node[below=0.3cm of node, font=\scriptsize\itshape, text=gray]
  {machine fingerprint};

\end{tikzpicture}
\caption{Customer signup and onboarding flow}
\end{figure}

\subsection{Login Flow}

Two separate login mechanisms serve two different contexts:

\begin{table}[H]
\centering
\renewcommand{\arraystretch}{1.4}
\begin{tabularx}{\textwidth}{@{}lXX@{}}
\toprule
 & \textbf{Portal login} & \textbf{Dashboard login} \\
\midrule
\textbf{Server}    & License/Control Server (hosted)    & Platform Server (local node) \\
\textbf{Credential}& Email + password + optional TOTP   & API key only \\
\textbf{Token}     & JWT (24h expiry, refreshable)       & API key (long-lived, revocable) \\
\textbf{Access}    & Account management, billing, keys  & Agents, sessions, memory, CLI \\
\bottomrule
\end{tabularx}
\caption{Portal login vs.\ local dashboard login — different credentials, different contexts}
\end{table}

\section{API Key Management}

API keys are the primary authentication mechanism for programmatic access to
the Platform Server. They are issued by the License/Control Server and
validated locally by the Platform Server — no round-trip to the hosted server
is required for every request.

\subsection{Key Lifecycle}

\begin{figure}[H]
\centering
\begin{tikzpicture}[
  st/.style={rectangle, rounded corners=3pt, draw=CBiz!60, fill=CBiz!7,
             text width=2.6cm, align=center, minimum height=0.85cm,
             font=\footnotesize\bfseries},
  node distance=0.8cm
]
\node[st] (create) {Create in\\portal};
\node[st, right=of create] (display) {Displayed\\once};
\node[st, right=of display] (store)  {Stored as\\SHA-256 hash};
\node[st, right=of store]   (use)    {Used via\\X-API-Key header};
\node[st, right=of use]     (revoke) {Revoked in\\portal};

\foreach \a/\b in {create/display,display/store,store/use,use/revoke}
  \draw[arr=CBiz!60] (\a) -- (\b);

\node[below=0.3cm of display, font=\scriptsize\itshape, text=CWarn]
  {copy immediately};
\node[below=0.3cm of store, font=\scriptsize\itshape, text=gray, align=center]
  {never stored in\\plaintext};

\end{tikzpicture}
\caption{API key lifecycle — plaintext is shown only once; only the hash is persisted}
\end{figure}

\subsection{Key Structure}

\begin{defbox}{API Key Format}
\texttt{ck\_live\_\textless{}32-char-random-bytes\textgreater{}}

\smallskip
\begin{itemize}[noitemsep]
  \item \texttt{ck\_} — Connector key namespace
  \item \texttt{live\_} or \texttt{test\_} — Environment discriminator
  \item 32 random bytes (Base58 encoded) — 192 bits of entropy
  \item \textbf{Prefix \texttt{ck\_live\_a1b2c3d4}} visible for identification
  \item \textbf{Full key stored as SHA-256 hash only}
\end{itemize}
\end{defbox}

Keys carry a \texttt{scopes} list that limits what operations they can perform.
Scope tokens follow the pattern \texttt{<noun>:<action>}, e.g.\
\texttt{agents:read}, \texttt{agents:write}, \texttt{decisions:write}.

\section{License Activation}

\subsection{Activation Flow}

After purchasing a plan, the customer runs a single command on their node:

\begin{lstlisting}[language=bash]
connectorctl activate <license-key>
\end{lstlisting}

The node contacts the License/Control Server, which verifies the key, generates
a machine fingerprint binding, and returns a signed license file. The license
file is stored at \texttt{/etc/connector/license.key} and read at every boot.

\subsection{License Tiers}

The codebase currently contains more than one entitlement model. For business,
sales, and documentation purposes, Connector should present one
\textbf{canonical 2026 packaging model}: a small self-serve base, a team plan,
an enterprise plan, and a paid pilot overlay used to land strategic accounts.

\begin{table}[H]
\centering
\renewcommand{\arraystretch}{1.5}
\rowcolors{2}{gray!5}{white}
\begin{tabularx}{\textwidth}{@{}lXXXX@{}}
\toprule
\textbf{Dimension} & \textbf{Community} & \textbf{Pro} & \textbf{Team} & \textbf{Enterprise} \\
\midrule
Target buyer & Individual builder & Small production team & Cross-functional AI team & Regulated or large enterprise \\
Monthly price & Free & \$49--\$99 & \$299--\$999 & Custom / annual contract \\
Agents & 2--3 & Up to 20 & High or uncapped & Contracted \\
Tokens & Small daily/monthly cap & 500K monthly class & 5M monthly class & Contracted or uncapped \\
Memory retention & 30 days & 1 year & Unlimited & Unlimited \\
HIPAA / advanced compliance & -- & -- & Optional add-on / controlled release & Included by contract \\
SSO / advanced admin & -- & -- & Optional & Included \\
On-prem license & -- & -- & Limited cases & Core offer \\
Support & Community & Email & Priority & Dedicated \\
\bottomrule
\end{tabularx}
\caption{Recommended canonical packaging for Connector in 2026}
\end{table}

\subsection{Paid pilot overlay}

Paid pilots should not be treated as ad hoc discounting. Internally, the
control server already models pilots as \texttt{PilotGrant} overlays applied on
top of a base entitlement. Commercially, this is the right design because it
avoids creating one-off tiers for every prospect.

\begin{table}[H]
\centering
\renewcommand{\arraystretch}{1.45}
\begin{tabularx}{\textwidth}{@{}lX@{}}
\toprule
\textbf{Pilot field} & \textbf{Commercial meaning} \\
\midrule
\texttt{tier\_override} & Temporarily grant a higher plan for evaluation \\
\texttt{agent\_limit\_override} & Expand concurrency for one team or workflow \\
\texttt{packet\_limit\_override} & Expand memory / evidence retention for the pilot \\
\texttt{features\_override} & Turn on gated features such as HIPAA, SSO, or multi-cell \\
\texttt{expires\_at} & Force a commercial decision point instead of an endless trial \\
\bottomrule
\end{tabularx}
\caption{Pilot overlays — commercial meaning of the internal pilot grant model}
\end{table}

All license tiers and pilot overlays produce a signed \texttt{License} object
with an Ed25519 signature. The Platform Server verifies this signature at boot
and on every agent start without a network round-trip.

\section{Billing Integration}

Connector uses \textbf{Stripe} for all payment processing. The billing flow
is fully event-driven through webhooks — the License/Control Server reacts
to Stripe events to update entitlements in real time.

\begin{figure}[H]
\centering
\begin{tikzpicture}[
  box/.style={rectangle, rounded corners=4pt, draw=#1!70, fill=#1!8,
              text width=3cm, align=center, minimum height=1.1cm,
              font=\small\bfseries, inner sep=5pt},
  node distance=1.2cm
]
\node[box=CBiz] (cust)  {Customer\\selects plan};
\node[box=CBiz, right=of cust] (stripe) {Stripe\\Checkout};
\node[box=CBiz, right=of stripe] (wh)   {Webhook\\handler};
\node[box=CBiz, right=of wh]   (upd)   {Update\\entitlements};
\node[box=CKernel, right=of upd] (key)  {Issue\\license key};

\foreach \a/\b in {cust/stripe,stripe/wh,wh/upd,upd/key}
  \draw[arr=CBiz!60] (\a) -- (\b);

\node[below=0.25cm of stripe, font=\scriptsize\itshape, text=gray, align=center]
  {Stripe-hosted\\form};
\node[below=0.25cm of wh, font=\scriptsize\itshape, text=gray, align=center]
  {\texttt{subscription.created}\\verified via HMAC};

\end{tikzpicture}
\caption{Stripe billing flow — webhook-driven entitlement updates}
\end{figure}

\subsection{Handled Stripe Events}

\begin{table}[H]
\centering
\renewcommand{\arraystretch}{1.4}
\begin{tabularx}{\textwidth}{@{}lX@{}}
\toprule
\textbf{Stripe event} & \textbf{Connector action} \\
\midrule
\texttt{customer.subscription.created}  & Upgrade tier, issue license key \\
\texttt{customer.subscription.updated}  & Adjust agent/memory/token limits \\
\texttt{customer.subscription.deleted}  & Downgrade to Free; agents above limit paused \\
\texttt{invoice.payment\_failed}        & Notify customer; 7-day grace period \\
\texttt{invoice.payment\_succeeded}     & Clear any payment-failure flag \\
\bottomrule
\end{tabularx}
\caption{Stripe webhook events and their Connector side effects}
\end{table}

\section{Usage Tracking and Telemetry}

The Platform Server reports a minimal telemetry payload to the
License/Control Server every 15 minutes. This payload contains only
aggregate counters — no agent data, no MemPacket content, no user queries.

\begin{table}[H]
\centering
\renewcommand{\arraystretch}{1.4}
\begin{tabularx}{\textwidth}{@{}lXl@{}}
\toprule
\textbf{Metric} & \textbf{Description} & \textbf{Used for} \\
\midrule
\texttt{agents\_running}     & Active agent count at time of report & Limit enforcement \\
\texttt{memory\_packets}     & Total MemPacket count                & Storage billing \\
\texttt{tokens\_consumed}    & Cumulative LLM tokens this period    & Token budget billing \\
\texttt{uptime\_seconds}     & Node uptime                          & SLA calculation \\
\texttt{session\_count}      & Sessions opened this period          & Usage graphs \\
\texttt{tool\_invocations}   & External tool calls this period      & Usage graphs \\
\bottomrule
\end{tabularx}
\caption{Telemetry metrics reported to the License/Control Server — aggregates only}
\end{table}

These metrics power the customer dashboard usage graphs, admin-panel
fleet analytics, and (in future releases) usage-based pricing for token
consumption above tier limits.

\section{Go-to-Market Motion}

Connector should use a \textbf{hybrid bottom-up + top-down motion}. The initial
user is usually technical: a platform engineer, AI engineer, or Head of AI who
needs to explain or control one production workflow. The deal closes when that
technical win becomes legible to a broader stakeholder set: compliance,
security, and the executive sponsor responsible for rollout risk.

\begin{table}[H]
\centering
\renewcommand{\arraystretch}{1.45}
\begin{tabularx}{\textwidth}{@{}lXX@{}}
\toprule
\textbf{Motion layer} & \textbf{Primary goal} & \textbf{Recommended execution} \\
\midrule
Bottom-up & Create urgency and technical conviction & CLI demo, local node install, one broken workflow explained end-to-end \\
Founder-led sales & Convert interest into a paid pilot & Direct outreach to 30--50 target accounts in healthcare, fintech, and enterprise AI platform teams \\
Top-down validation & Remove deal blockers & Bring in CTO, compliance, and security reviewers early with architecture and data-sovereignty materials \\
Land-and-expand & Grow from one workflow to one business unit & Add quotas, policy control, audit evidence, and more workloads after pilot success \\
\bottomrule
\end{tabularx}
\caption{Recommended GTM motion for Connector}
\end{table}

\subsection{Why this motion fits the product}

Menlo Ventures' 2025 enterprise AI findings matter here: AI products are being
purchased more often than built, they move to production at unusually high
rates once evaluation starts, and a meaningful share of AI spend is initiated
bottom-up. Connector should therefore \textbf{land technically and close
organizationally}.

\section{2026 Entry Plan}

\subsection{Primary commercial objective}

The 2026 objective should be to prove that Connector is not merely a compelling
architecture, but a repeatable commercial wedge. A practical target is:

\begin{conceptbox}{2026 core target}
\centering
\textbf{10 paid pilots, each priced at roughly \$5k--\$7k MRR equivalent,}\\[2pt]
\textbf{with at least 4 pilot-to-annual conversions by year end.}
\end{conceptbox}

In practice, that pilot price can be packaged as a \textbf{90-day paid pilot at
\$15k--\$21k total}, which preserves the MRR-equivalent narrative while fitting
enterprise procurement more naturally than a pure month-to-month experimental
contract.

\subsection{Recommended target account profile}

The first 10 pilots should not be spread across every possible industry. They
should be concentrated where Connector's wedge is strongest:

\begin{itemize}[leftmargin=1.5em,itemsep=2pt]
  \item \textbf{4 healthcare / healthtech pilots} for patient operations,
        triage, prior authorization, or compliance-heavy workflows.
  \item \textbf{3 fintech / insurtech pilots} where decisions must be explainable
        to internal auditors or customer support teams.
  \item \textbf{3 enterprise platform pilots} where internal AI platform teams are
        trying to standardize agent operations across multiple teams.
\end{itemize}

This focus improves message discipline, sales efficiency, and referenceability.

\subsection{Pilot structure}

\begin{table}[H]
\centering
\renewcommand{\arraystretch}{1.45}
\begin{tabularx}{\textwidth}{@{}lX@{}}
\toprule
\textbf{Pilot element} & \textbf{Recommended structure} \\
\midrule
Duration & 8--12 weeks, with explicit success criteria signed before kickoff \\
Commercial form & \$15k--\$21k fixed pilot fee (equivalent to \$5k--\$7k MRR) \\
Scope & One workflow, one team, one production-like environment \\
Success criteria & Show one failure chain, one policy decision path, one signed evidence trail, one operator workflow improvement \\
Technical deliverable & Connector node deployment, portal account, API keys, pilot overlay entitlement, usage dashboard \\
Conversion path & Annual Team or Enterprise contract with expansion to more workflows or business units \\
\bottomrule
\end{tabularx}
\caption{Recommended paid pilot design}
\end{table}

\subsection{Quarter-by-quarter execution plan}

\begin{table}[H]
\centering
\renewcommand{\arraystretch}{1.45}
\begin{tabularx}{\textwidth}{@{}lX@{}}
\toprule
\textbf{Period} & \textbf{Execution objective} \\
\midrule
Q1 2026 & Finalize narrative, tighten demo, ship stable onboarding, identify 30--50 target accounts, close first 2 paid pilots \\
Q2 2026 & Reach 4--5 paid pilots, publish 2 detailed case studies or design-partner stories, refine packaging and pricing based on objections \\
Q3 2026 & Reach 7--8 paid pilots, convert first 2--3 pilots to annual contracts, formalize repeatable founder-led sales process \\
Q4 2026 & Reach 10 paid pilots total, convert at least 4 to annual contracts, use the data room to support pre-seed or seed fundraising \\
\bottomrule
\end{tabularx}
\caption{2026 execution plan — from wedge validation to repeatable GTM}
\end{table}

\section{Fundraising Narrative}

\subsection{What the pre-seed story should be}

The pre-seed story is not ``we built a very deep AI operating system.'' The
pre-seed story is:

\begin{warnbox}{Fundraising message discipline}
Connector helps enterprises move AI into production by giving them a governed
runtime and evidence layer. The company has a clear entry wedge, paid pilot
traction, strong expansion logic, and a product architecture aligned with data
sovereignty requirements.
\end{warnbox}

Founder Institute's 2025 benchmark guidance is useful here: in 2026 investors
want paid traction, capital efficiency, and a concrete plan for what the next
round of capital unlocks. For Connector, that means the fundraise should be
anchored in \textbf{evidence of demand plus clarity of use of proceeds}.

\subsection{Recommended pre-seed milestones}

\begin{table}[H]
\centering
\renewcommand{\arraystretch}{1.45}
\begin{tabularx}{\textwidth}{@{}lX@{}}
\toprule
\textbf{Milestone area} & \textbf{What investors should see} \\
\midrule
Commercial traction & 4--6 active paid pilots minimum; 10 is a strong year-end target \\
Revenue quality & Clear pilot pricing discipline and visible path to annual contracts \\
Category clarity & One-sentence wedge buyers immediately understand \\
Referenceability & At least 2 strong customer stories in regulated or high-consequence workflows \\
Capital plan & Hiring and roadmap tied to specific go-to-market and product bottlenecks \\
\bottomrule
\end{tabularx}
\caption{Milestones that support a credible pre-seed raise}
\end{table}

\subsection{Use of proceeds}

If Connector raises a pre-seed round, the capital should primarily fund:

\begin{itemize}[leftmargin=1.5em,itemsep=2pt]
  \item \textbf{Product hardening} — onboarding, pilot tooling, entitlement
        consistency, dashboards, packaging polish.
  \item \textbf{Customer success and implementation} — fast pilot launches and
        clean conversions.
  \item \textbf{Founder-led GTM leverage} — tighter outbound process, repeatable
        demo assets, sector-specific collateral.
  \item \textbf{Compliance readiness} — the materials and controls required for
        enterprise trust in healthcare and fintech sales cycles.
\end{itemize}

\section{Commercial Scorecard}

The business should be managed against a small number of metrics rather than a
large vanity dashboard.

\begin{table}[H]
\centering
\renewcommand{\arraystretch}{1.45}
\begin{tabularx}{\textwidth}{@{}lXX@{}}
\toprule
\textbf{Metric} & \textbf{Why it matters} & \textbf{2026 target} \\
\midrule
Paid pilots closed & Proves the wedge is commercially real & 10 \\
Pilot MRR-equivalent & Shows pricing discipline and revenue quality & \$50k--\$70k combined \\
Pilot-to-annual conversion & Shows the product expands beyond curiosity & 40--50\% \\
Average pilot launch time & Shows operational readiness & Under 30 days from signature \\
Champion-to-exec progression & Shows multi-stakeholder sale health & In every paid pilot \\
Gross logo churn & Protects early narrative quality & Near zero during pilot year \\
\bottomrule
\end{tabularx}
\caption{Commercial scorecard for the first year of GTM}
\end{table}

\section{SWOT Analysis}

\begin{table}[H]
\centering
\renewcommand{\arraystretch}{1.45}
\begin{tabularx}{\textwidth}{@{}lX@{}}
\toprule
\textbf{Dimension} & \textbf{Connector assessment} \\
\midrule
Strengths & Strong technical differentiation, data-sovereign architecture, runtime-native governance, explainability with proof, clear value in regulated workflows \\
Weaknesses & Broad architecture can be mis-pitched, category education is still required, enterprise trust must be earned, packaging needs tight canonicalization \\
Opportunities & Enterprise AI budgets are growing, buyers increasingly purchase instead of build, agent governance is still early, regulated verticals need practical auditability \\
Threats & Larger incumbents may add partial governance features, category confusion may slow deals, long enterprise cycles can starve a small team, broad platform positioning can dilute the wedge \\
\bottomrule
\end{tabularx}
\caption{SWOT summary for Connector's 2026 business posture}
\end{table}

\begin{insightbox}{One sentence to keep the company honest}
Connector should not try to replace the world before it earns trust in one
workflow. The product wins when it proves one broken AI decision path,
controls it, and then expands from that foothold.
\end{insightbox}

%%────────────────────────────────────────────────────────────
\chapter*{Ground Truth: Pre-Revenue Founder Economics}
\addcontentsline{toc}{chapter}{Ground Truth: Pre-Revenue Founder Economics}
%%────────────────────────────────────────────────────────────

The preceding sections of this chapter describe where the business should be.
This chapter describes how two founders with no revenue actually get there.
Every number here is drawn from real startup data, not investor storytelling.

\begin{warnbox}{The salary problem is a business problem}
A founder who cannot pay rent cannot execute. This is not a personal
weakness --- it is a planning failure. If you go full-time without solving
the salary question first, you will make desperate decisions (wrong customers,
wrong pricing, wrong hires) within four months. Solve it in writing before day
one of full-time work.
\end{warnbox}

\section*{Founder Runway: How to Pay Two Salaries Before Any Revenue}

The minimum monthly requirement for two founders to survive full-time is
honest and non-negotiable:

\begin{table}[H]
\centering
\renewcommand{\arraystretch}{1.45}
\begin{tabularx}{\textwidth}{@{}lXl@{}}
\toprule
\textbf{Item} & \textbf{Details} & \textbf{Monthly cost} \\
\midrule
Founder 1 living wage & Rent, food, transport, health cover (modest) & \$3,500--5,000 \\
Founder 2 living wage & Same basis & \$3,500--5,000 \\
Tooling \& infrastructure & GitHub, cloud credits, domain, legal templates & \$400--600 \\
Travel for pilot meetings & 1--2 in-person visits per quarter & \$500--800 \\
\midrule
\textbf{Total monthly burn} & \textbf{Pre-revenue, two founders} & \textbf{\$8,000--11,000} \\
\bottomrule
\end{tabularx}
\caption{Realistic monthly burn for two founders before any revenue}
\end{table}

\noindent Four mechanisms fund this. Choose one or combine them:

\begin{enumerate}[leftmargin=1.8em, itemsep=4pt]
  \item \textbf{Pre-seed SAFE from angels (\$75K--\$150K).}\;
    One or two angel investors who understand developer-infrastructure
    businesses. At a \$2M--\$3M cap SAFE, \$100K gives 12--15 months of
    founder salaries at minimal living rates. This is the cleanest option
    --- no board seat, no preferred shares, fast to close. Angels at this
    stage are usually former founders, CTOs, or senior engineers at AI
    companies. They invest in people and in conviction about the problem,
    not in a product that works.

  \item \textbf{One founder keeps a part-time consulting income.}\;
    One day per week of senior engineering consulting (\$1,200--2,000/day)
    = \$5,000--8,000/month from one founder. This extends runway indefinitely
    but costs roughly 20\% of that founder's product time. Use this as a
    bridge, not a permanent state. The ceiling is 3--4 months before it
    noticeably slows the company.

  \item \textbf{First pilot revenue within 90 days.}\;
    Two pilot customers at \$2,000/month each = \$4,000 MRR. Combined with
    minimal savings, this covers both founder salaries. This is achievable
    but requires the founders to have built a working demo and a short list
    of warm contacts before going full-time. Do not go full-time and
    \emph{then} start building the demo.

  \item \textbf{Government and accelerator non-dilutive grants.}\;
    In the UK (Innovate UK), EU (EIC Accelerator), Canada (SR\&ED), and
    USA (SBIR Phase I for AI infrastructure), grants of \$25K--\$150K are
    available for AI systems research with a novel component. Connector's
    predicate-chain governance and CLS compiler are genuinely novel.
    Timeline: 3--6 months to apply and receive. Not fast enough to be
    the primary bridge, but valuable as a secondary source.
\end{enumerate}

\begin{insightbox}{The honest decision tree}
If you have at least 4 months of personal savings each \emph{and} a working
demo \emph{and} a list of 20 warm contacts: go full-time and close pilot 1
within 90 days. If you do not have all three, line up the angel SAFE first.
Going full-time without runway is not hustle --- it is bad planning.
\end{insightbox}

\section*{Who Actually Pays During Pilots}

The commercial scorecard in the prior section targets \$50K--\$70K combined
pilot MRR. That number is real but the path to it does not go through
large enterprises.

\begin{warnbox}{Enterprise sales will not save 2026}
Fortune~500 companies, large banks, and hospital systems have procurement
cycles of 6--18~months. They require SOC2~Type~II (9--12~months and
\$30K--\$80K to obtain), a security questionnaire response that takes
2--3~weeks per account, an infosec review, a legal MSA, a DPA, and an SLA
with financial penalties. Chasing these accounts in 2026 means zero revenue
in 2026. Do not do it.
\end{warnbox}

\noindent The ideal pilot customer in 2026 looks like this:

\begin{table}[H]
\centering
\renewcommand{\arraystretch}{1.5}
\begin{tabularx}{\textwidth}{@{}lX@{}}
\toprule
\textbf{Dimension} & \textbf{Ideal profile} \\
\midrule
Company stage & Funded AI startup (post-seed to Series~A), or boutique AI
               consulting firm, or Series~A/B SaaS company actively adding
               AI agent features \\
Company size  & 8--60 employees. Engineering team of 5--25. \\
Funding raised & \$2M--\$20M. Enough budget to move without a committee, not
                 so large they need procurement. \\
Buying decision & CTO or VP Engineering makes the call alone or with one other
                 person. No legal committee, no CISO review in round 1. \\
Technical profile & Team already using LLMs in production. Have already felt
                   the pain of agents making decisions they cannot explain or
                   audit. \\
Compliance pressure & At least one of: HIPAA, SOC2 (in progress), GDPR, or
                      a regulated-industry client they are selling to. \\
\midrule
\textbf{Concrete examples} & Early digital health startup building AI clinical
  decision support. AI-native fintech startup with LLM-driven underwriting.
  Boutique AI consulting firm that needs governance as a differentiator for
  their own client proposals. Series~A developer-tools company adding
  AI features and worried about audit trails. \\
\bottomrule
\end{tabularx}
\caption{Ideal pilot customer profile for 2026 --- not enterprise, not consumer}
\end{table}

\subsection*{Why These Customers Pay}

The question is not just who pays --- it is why they pay when Connector is
early and the codebase is not yet fully polished.

\begin{enumerate}[leftmargin=1.8em, itemsep=5pt]
  \item \textbf{Build vs.\ buy arithmetic.}\;
    A funded startup burning \$200K/month on salaries would need 2--3
    senior engineers for 12--18 months to build what Connector delivers.
    That is \$300K--\$600K in engineering cost. A \$2,000--\$3,000/month
    pilot is not a purchase decision --- it is a cost-of-alternatives
    decision. The CTO's mental model is: ``could we build this ourselves?''
    When the answer is ``yes but it would take us 18 months,'' Connector
    wins.

  \item \textbf{Design partner access.}\;
    At the \$1,500--\$3,000/month pilot price, the customer is not buying
    finished software. They are buying direct access to the founders who
    will build exactly the feature they need next. Founders respond to their
    Slack messages on the same day. Enterprise vendors do not. Early
    customers pay for that relationship --- it is a real competitive advantage
    for them internally.

  \item \textbf{Competitive differentiation.}\;
    In regulated or compliance-adjacent markets, a startup that can say
    ``our AI agents produce a cryptographic proof of every decision'' to
    their own enterprise clients wins deals that their competitors lose.
    The pilot customer is not buying Connector for themselves --- they are
    buying it to use as a sales argument to \emph{their} enterprise customers.

  \item \textbf{Urgency.}\;
    Funded startups have a board. The board asks about AI governance now.
    The CTO has 6--8 weeks to have an answer. Connector is the answer they
    can deploy in days, not months. That urgency converts to a paid decision
    without a procurement process.
\end{enumerate}

\section*{10 Pilots, 2 People: The Actual Playbook}

Ten paid pilots from two founders in 12 months is achievable. It requires
a specific method, not general hustle.

\begin{warnbox}{The codebase condition is honest}
As of early 2026, the Connector platform has excellent architectural depth
(Memory Kernel, CLS compiler, CNP protocol, SOE) but lacks the operational
maturity that makes self-service onboarding possible: a one-command
deployment, complete API documentation, a stable demo environment, and
polished error messages. This means the first 5 pilots must be
\textbf{concierge-style} --- the founders do the setup, hold the customer's
hand, and manually resolve friction. This is not a weakness to hide. It
is the correct strategy: you learn 10$\times$ faster and customers feel
10$\times$ more valued.
\end{warnbox}

\begin{table}[H]
\centering
\renewcommand{\arraystretch}{1.5}
\begin{tabularx}{\textwidth}{@{}lllXl@{}}
\toprule
\textbf{Phase} & \textbf{Period} & \textbf{Target} & \textbf{Method} & \textbf{MRR} \\
\midrule
Phase 0 & Jan--Feb 2026 & Demo-ready & Fix deployment: Docker Compose
  single-node, 5-page integration doc, 3 worked examples, stable local demo
  that runs in under 10 minutes & \$0 \\
Phase 1 & Feb--Apr 2026 & 2 pilots & Direct outreach to 60--80 funded AI
  startups via LinkedIn, YC company list, Twitter/X. Founders do every
  demo personally. Concierge onboarding with weekly calls. Price:
  \$1,500--\$2,000/month & \$3K--\$4K \\
Phase 2 & May--Jul 2026 & 5--6 pilots & Use Phase~1 learnings to write a
  one-page onboarding guide. Outreach to 100--150 targets with a social
  proof statement (``2 companies using this in production''). SDK hire
  joins to absorb integration support & \$10K--\$15K \\
Phase 3 & Aug--Oct 2026 & 10 pilots & Semi-self-service. Referrals from
  Phase~1 customers. Sales hire joins to run outbound pipeline. Price
  increases to \$2,500--\$3,500/month as the product matures &
  \$22K--\$30K \\
\bottomrule
\end{tabularx}
\caption{Phased pilot capture plan --- 2 founders, concierge model, no enterprise}
\end{table}

\noindent The outreach mechanics are not complex: identify companies by
searching for ``LLM agents'', ``AI workflow'', or ``AI governance'' on LinkedIn
job postings, AngelList, and Crunchbase. Companies actively hiring for these
roles have the problem Connector solves. Message the CTO directly. The message
is two sentences: what you solve and who you have already helped. No deck.
No brochure. A demo link that runs in the browser in 5 minutes.

\section*{The Three Hires: Who, Why, and When}

\begin{warnbox}{Wrong hires destroy early startups faster than no hires}
An enterprise AE costs \$120K--\$180K/year OTE and expects a polished product,
a clean CRM, and a well-defined ICP. They will leave after 90 days if they
cannot close. An RTC engineer or distributed-systems researcher costs
\$160K--\$220K/year and will spend 6 months understanding the codebase before
shipping anything customer-visible. Neither is the right hire in 2026. The
right hires are people who can do multiple jobs badly rather than one job
perfectly.
\end{warnbox}

\begin{table}[H]
\centering
\renewcommand{\arraystretch}{1.7}
\begin{tabularx}{\textwidth}{@{}lXlll@{}}
\toprule
\textbf{Hire} & \textbf{Role and profile} & \textbf{Timing} &
\textbf{Monthly comp} & \textbf{Equity} \\
\midrule
Tech~1 & \textbf{Product engineer (Rust / full-stack).} Reads the existing
  codebase, ships customer-visible features, handles deployments at pilot
  sites, helps debug integration issues. Has worked at a startup before.
  5--8 years. Comfortable with ambiguity. \emph{Not} a research engineer ---
  you do not have a research problem yet, you have a delivery problem.
  & After pilot 3, MRR $\geq$\$6K & \$5,500--7,500/mo & 0.75--1.25\% \\
\midrule
Tech~2 & \textbf{Customer/SDK engineer (Python + TypeScript).} Builds the
  Python and TypeScript SDKs, writes integration examples, runs onboarding
  calls with the technical contact at each pilot. Part sales engineer, part
  customer success, part developer. Must be comfortable on video calls with
  strangers. 3--6 years. Background in API platforms or developer tools is
  ideal.
  & After pilot 5, MRR $\geq$\$12K & \$4,500--6,500/mo & 0.5--0.75\% \\
\midrule
Sales~1 & \textbf{Founder-minded BDR/AE (developer tools background).} Does
  outbound research, writes personalised outreach, follows up, runs demos
  alongside the founder, and closes \$2K--\$4K/month deals. Has worked at
  an early-stage API or developer-tools startup (10--50 employees) before.
  Is comfortable with a product that is still rough. Is not expecting a CRM
  and a polished pitch deck on day one. Motivated by equity upside, not just
  commission.
  & After pilot 6, MRR $\geq$\$16K & \$4,000--5,500/mo + deal bonus & 0.3--0.5\% \\
\bottomrule
\end{tabularx}
\caption{First three hires --- timing, profile, and honest compensation for 2026}
\end{table}

\subsection*{Why Not Hire Earlier}

Hiring before revenue reaches the thresholds above means the company is
paying salaries that the pilots do not yet cover. That shortens the runway
without accelerating the outcome. The two founders can operate:

\begin{itemize}[leftmargin=1.5em, itemsep=3pt]
  \item Pilot 1--2: Founder A owns product and architecture. Founder B owns
        outreach, demos, and onboarding. Both do customer calls.
  \item Pilot 3--5: One founder is now spending 60\% of time on customer
        integration issues. That is the signal to hire Tech~1.
  \item Pilot 5--7: Integration issues are resolved by Tech~1 but SDK
        surface and onboarding repeatability is still fragile. That is the
        signal to hire Tech~2.
  \item Pilot 7--9: Outbound is now the bottleneck. The founder cannot run
        sales, manage three hires, and keep the product moving. That is the
        signal to hire Sales~1.
\end{itemize}

\section*{Honest Monthly Cash-Flow Projection}

\begin{table}[H]
\centering
\renewcommand{\arraystretch}{1.5}
\begin{tabularx}{\textwidth}{@{}lXXXX@{}}
\toprule
\textbf{Month} & \textbf{Pilots} & \textbf{MRR} & \textbf{Monthly burn} &
\textbf{Net monthly} \\
\midrule
1--3  & 0 & \$0 & \$9K (founders only) &
         $-$\$9K (funded by pre-seed or savings) \\
4--5  & 2 & \$4K & \$9K & $-$\$5K (still drawing down) \\
6--7  & 4 & \$9K & \$9K & $\pm$\$0 (cash-flow neutral) \\
8--9  & 6 + Tech~1 joins & \$15K & \$15K & $\pm$\$0 \\
10--11 & 8 + Tech~2 joins & \$21K & \$20K & $+$\$1K \\
12    & 10 + Sales~1 joins & \$27K--\$30K & \$24K &
         $+$\$3K--\$6K cash-flow positive \\
\bottomrule
\end{tabularx}
\caption{Honest 12-month cash-flow model --- two founders, no external capital after pre-seed}
\end{table}

\noindent At month 12 with 10 pilots at an average of \$2,700/month, Connector
is at \$324K ARR with three hires and two founders generating a small positive
monthly cash flow. This is not an ambitious projection --- it is the
\textbf{minimum} viable outcome from the 2026 plan. It funds the story for
the Seed~1 round (\$1M--\$2M) at a valuation that reflects real revenue
traction, not pre-revenue promise.

\begin{insightbox}{The goal of 2026 is not \$1M ARR}
The goal of 2026 is to prove that a specific customer pays a specific price
for a specific outcome --- reliably, repeatedly, without the founders being
in every call. That proof is worth more to a Seed investor than any amount
of unpaid pilots or letters of intent. Eleven paying customers at
\$2,500/month beats fifty LOIs every time.
\end{insightbox}

\section*{What ``Pilot-Ready'' Means for This Codebase}

The Connector codebase has architectural depth that most early-stage
infrastructure companies take two additional years to reach. The Memory
Kernel, CLS compiler, CNP protocol, SOE, and cognitive substrate are
genuine technical assets. However, the gap between ``architecturally
correct'' and ``a customer can use this next Tuesday'' is significant and
must be closed before outreach begins.

\begin{table}[H]
\centering
\renewcommand{\arraystretch}{1.5}
\begin{tabularx}{\textwidth}{@{}lXll@{}}
\toprule
\textbf{Component} & \textbf{Current state} & \textbf{Gap to pilot-ready} &
\textbf{Effort} \\
\midrule
Memory Kernel      & Core implemented, CID working & Stable public API surface,
                     5 documented operations & 2--3 weeks \\
Deployment         & Requires manual setup, multiple steps & Docker Compose
                     single-node that works in one command & 1--2 weeks \\
Demo environment   & Internal, requires context to run & Browser-accessible
                     demo with a scripted walkthrough & 2 weeks \\
Python SDK         & Not yet customer-facing & 10-endpoint coverage, pip
                     installable, 3 worked examples & 3--4 weeks \\
API documentation  & Internal only & 15-page integration guide covering
                     the 5 pilot workflows & 2 weeks \\
Error messages     & Technical / internal & Human-readable, actionable, no
                     internal stack traces to customers & 1 week \\
CLS compiler       & Design complete, in development & A working compiler
                     for 3 contract patterns used in pilot workflows & 4--6 weeks \\
Billing \& licensing & Architecture built & End-to-end test: sign up, activate,
                     use, invoice & 1--2 weeks \\
\midrule
\textbf{Total} & & & \textbf{8--12 weeks, 2 founders full-time} \\
\bottomrule
\end{tabularx}
\caption{Gap analysis: current codebase state to pilot-ready}
\end{table}

\noindent This 8--12 week window (January--March 2026) is \textbf{Phase~0}.
No outreach happens during this phase. No demos. No sales conversations.
The two founders build the demo environment, write the first integration
guide, fix deployment, and get the Python SDK to a state where a CTO can
try it in an afternoon without calling anyone. Only then does Phase~1
outreach begin.

\begin{warnbox}{Do not start selling before the demo works}
Reaching out to 80 potential pilot customers before the demo runs reliably
wastes the only asset you cannot recover: warm introductions. A CTO who
tries the demo, hits an error, and does not get a response within 2 hours
will never be a customer. You get one first impression per person.
Use Phase~0 to make it count.
\end{warnbox}

%%────────────────────────────────────────────────────────────
\section*{Fundraising Reality: The Honest Timeline}
%%────────────────────────────────────────────────────────────

The fundraising timeline in the growth chapter requires correction against
execution reality. Raising too early produces worse terms, wasted founder
time, and investor relationships built on a story that has no proof yet.

\begin{warnbox}{Do not raise a pre-seed in 2026}
Going to angels with zero paying customers produces a \$1--1.5M cap SAFE
at best. Going to the same angels in Q4~2026 with \$15K--25K MRR and
8 paying customers produces a \$2.5--3.5M cap SAFE. That difference
compounds forever: it means \emph{less dilution for the same cash}.
The right strategy is to bootstrap through 2026, accumulate traction,
and raise on terms that reflect proof rather than promise.
\end{warnbox}

\begin{table}[H]
\centering
\renewcommand{\arraystretch}{1.6}
\begin{tabularx}{\textwidth}{@{}llXll@{}}
\toprule
\textbf{Round} & \textbf{Timing} & \textbf{Pre-conditions required} &
\textbf{Target raise} & \textbf{SAFE cap} \\
\midrule
Pre-seed & Q4~2026--Q1~2027 &
  8--10 paying pilots. \$15K--25K MRR. Founders paid from pilot
  revenue. Working demo. Clear ICP proven in two verticals. &
  \$75K--125K & \$2.5M--3.5M \\
\midrule
Seed~1 & Q4~2027 &
  15--20 paying customers. \$25K--40K MRR (\$300K--480K ARR). NRR above
  100\%. One reference customer willing to speak to investors. Team of
  2 founders plus 2--3 hires. ICP is narrow and proven, not broad. &
  \$500K--1.5M & \$4M--6M post-money \\
\midrule
Series~A & Q3--Q4~2028 &
  25--35 paying customers. \$600K--1M ARR. NRR above 110\% (customers
  expand, not just renew). SOC2~Type~II obtained. At least 3 customers
  paying above \$5K/month. Team of 8--12. Clear evidence of why Connector
  wins in its wedge. &
  \$2M--5M & \$12M--20M post-money \\
\midrule
Self-sufficient & End~2029 &
  \$1.5M--2M ARR. Monthly revenue exceeds monthly burn without raising
  further capital. Cash-flow positive or within 3 months of it. Team of
  12--18. Series~A capital still partially undeployed. &
  No raise needed & --- \\
\bottomrule
\end{tabularx}
\caption{Honest fundraising timeline with required pre-conditions}
\end{table}

\subsection*{Pre-Seed Reality: Why \$15--25K Angel Checks Are Hard}

The pre-seed round consists of 3--5 individual angel checks of \$15K--25K
each (total: \$75K--125K). On paper this sounds simple. In practice it is
friction-heavy:

\begin{itemize}[leftmargin=1.5em, itemsep=4pt]
  \item Each angel requires a 30--60 minute call, a 5-slide company summary,
        a SAFE term sheet, and follow-up. With 10--15 conversations to close
        4 cheques, the process takes 8--12 weeks of active founder time.
        During that period, one founder must keep running sales and product.

  \item Angels at this size often ask for MFN (Most Favoured Nation)
        provisions, pro-rata rights, or board observer seats. Agree only to
        pro-rata. The others add administrative overhead that kills momentum.

  \item The right angels are former founders or senior engineers at AI
        infrastructure companies (Databricks, Weights~\&~Biases, Scale~AI,
        Anthropic alumni). They invest in the \emph{problem} conviction, not
        the product state. Reaching them requires warm introductions, not
        cold emails.

  \item Rule: if you spend more than 90 days trying to close the pre-seed,
        stop and return to pilot execution. Revenue is a better bridge than
        a stalled fundraise.
\end{itemize}

\subsection*{Seed~1: What Investors Actually Need to See}

A Seed investor in 2027 is not funding a demo. They are funding a pattern:
a repeatable process that converts a specific type of customer into a paying
contract. The proof they require is not \emph{that} customers pay --- it is
\emph{why} they renew.

At the Seed~1 conversation, the founding team must be able to say:

\begin{itemize}[leftmargin=1.5em, itemsep=3pt]
  \item ``We have 15 customers. 12 have renewed past the initial 6-month
        pilot. 3 have expanded their contract.''
  \item ``Our average pilot-to-annual conversion is 60\%. Our NRR is 112\%.''
  \item ``Here is the exact workflow where we win: a Series~A AI startup or
        mid-market SaaS company that is deploying LLM agents in a
        compliance-sensitive context and needs cryptographic audit trails.''
  \item ``Here is the CTO who will take your call and explain exactly why
        they paid.''
\end{itemize}

\noindent Without these four statements grounded in real data, no credible
Seed fund writes a cheque. This is why the 2026 pilot execution is not a
``revenue target'' --- it is a \textbf{fundraising instrument}.

\subsection*{From Series~A to Self-Sufficient by End of 2029}

The Series~A funds the scaling of a proven motion. It is not for experiments.
The companies that fail between Seed and Series~A spend the Seed money
rebuilding the product for a new ICP rather than scaling the one that worked.

\begin{table}[H]
\centering
\renewcommand{\arraystretch}{1.5}
\begin{tabularx}{\textwidth}{@{}lXl@{}}
\toprule
\textbf{Quarter} & \textbf{What the Series~A funds} & \textbf{ARR target} \\
\midrule
Q1 2029 & 2 more sales hires (AE + SDR). First dedicated CS hire.
          SOC2 audit complete. & \$600K--800K \\
Q2 2029 & Enterprise motion starts: first 3 contracts at
          \$10K--15K/month. Product hardening for multi-tenant. & \$800K--1M \\
Q3 2029 & NRR driven by expansion: existing customers grow from
          \$3K to \$6K/month. New logo pipeline runs without founder. & \$1M--1.4M \\
Q4 2029 & Monthly cash-flow positive. Headcount 12--18. Founder salary
          at market rate (\$8K--12K/month each). No further capital required
          to keep the company alive. & \$1.5M--2M \\
\bottomrule
\end{tabularx}
\caption{Series A deployment and path to self-sufficiency by end of 2029}
\end{table}

\begin{insightbox}{Self-sufficient does not mean you stop raising}
``Self-sufficient by end of 2029'' means the company does not \emph{need}
to raise to survive. It may still choose to raise a Series~B to accelerate.
The difference is critical: raising from necessity means accepting any terms.
Raising from strength means choosing the right partner on your own schedule.
\end{insightbox}

%%────────────────────────────────────────────────────────────
\section*{Why Companies Like Connector Fail}
%%────────────────────────────────────────────────────────────

This section exists because most infrastructure startups with strong
technical architecture do not make it to \$1M ARR. The reasons are specific,
recurring, and entirely avoidable if named honestly in advance.

\begin{enumerate}[leftmargin=1.8em, itemsep=6pt]

  \item \textbf{ICP sprawl.}\;
    Connector's architecture can serve healthcare, fintech, education, and
    enterprise. If the team tries to sell to all four simultaneously, the
    product becomes generic and the sales message becomes vague. The
    customer says ``interesting'' and does not sign. The fix is to serve
    one ICP until you have 8 paying logos in that vertical, then expand.
    Connector's 2026 ICP is narrow: a funded AI startup or mid-market SaaS
    company deploying LLM agents with compliance pressure. Nothing else.

  \item \textbf{Architecture over delivery.}\;
    The Connector codebase has a Memory Kernel, a 6-pass CLS compiler, a
    cognitive substrate with 11 layers, and a 7-layer CNP protocol. None
    of these are interesting to a paying customer. What is interesting is
    a working audit trail for their specific agent workflow. Companies with
    deep architecture die when the team spends all year refining internals
    instead of shipping the one feature a customer is waiting for.

  \item \textbf{Wrong first customer.}\;
    One large enterprise excited about the platform in Q1 2026 is a trap.
    Their procurement process will consume 8 months. They will ask for
    SOC2, legal review, a custom MSA, and a sandbox environment that does
    not exist yet. By Q3 2026 the deal is dead, the team is demoralized,
    and they have no paying customers and no story for investors.

  \item \textbf{Founder cash starvation.}\;
    A founder who cannot pay rent makes three irrational decisions within
    six months: takes an investor's bad terms because they need the cash,
    accepts a customer's wrong requirements because they need the revenue,
    or burns out and walks away. This is the single most common cause of
    death for early infrastructure startups. It is solved in planning, not
    in execution.

  \item \textbf{Underpricing.}\;
    Pricing at \$500/month to be ``accessible'' means 100 customers to
    reach \$600K ARR. Two founders cannot manage 100 customer relationships
    while building the product. Pricing at \$3,000/month means 17 customers.
    The temptation to underprice is strong when rejection stings. The
    correct answer to rejection is not lowering the price --- it is finding
    a different customer who feels the pain more acutely.

  \item \textbf{Over-engineering the demo.}\;
    A technical demo that takes 30 minutes to explain and requires the
    founder to be present loses 60\% of potential customers before the
    value lands. The demo must answer one question in under 5 minutes:
    ``what breaks for you today that this fixes?'' Everything else is noise.

  \item \textbf{Chasing investors before customers.}\;
    Spending months in investor meetings without paying customers produces
    the same answer from every credible investor: ``come back when you have
    revenue.'' The time cost is 3--4 months of lost execution. The
    psychological cost is a demoralized team that feels like they are
    failing because investors said no. Investors do not validate the
    business. Customers do.

  \item \textbf{Single-founder technical dependency.}\;
    If only one person understands the codebase deeply and that person
    burns out, gets sick, or leaves, the company has a 30-day survival
    window. The Tech~1 hire must be someone who can eventually own the
    architecture independently. Documentation of internal systems is not
    optional --- it is company-survival infrastructure.

  \item \textbf{Free pilots.}\;
    Offering a free pilot ``to build goodwill'' achieves the opposite. It
    teaches customers that Connector has no pricing conviction, which signals
    that the product is not actually worth paying for. A customer who does
    not pay does not take onboarding seriously, does not give useful
    feedback, and cannot become a reference. Every pilot is paid from day
    one, even if the price is \$500/month. Price signals value.

  \item \textbf{Ignoring NRR until it is too late.}\;
    Acquiring 10 pilots and losing 6 is catastrophically worse than
    acquiring 6 and keeping all 6. A churn rate above 20\% in the first year
    is terminal --- it means the product is not solving a real problem or
    the customer was wrong to begin with. NRR must be reviewed after
    every pilot's 3-month mark. One churned customer is a signal. Two is
    a pattern that requires immediate product or ICP correction.

\end{enumerate}

\begin{warnbox}{The company can fail. This is how.}
Connector fails if: (a) both founders cannot pay themselves for more than
3 months and one of them takes a full-time job, (b) the first 3 pilot
attempts fail and the team pivots ICP without understanding why they
failed, or (c) the product ships a critical data breach or compliance
failure at a pilot customer before SOC2 is obtained. None of these are
technical failures. All three are planning failures. They are named here
so they can be planned against explicitly.
\end{warnbox}

%%────────────────────────────────────────────────────────────
\section*{What Pilots Actually Get: The Value, Not The Price}
%%────────────────────────────────────────────────────────────

The question that closes a pilot is not ``how much does it cost?'' It is
``what do I have at the end of six months that I do not have today?'' Every
pilot conversation starts with the deliverable, not the price. The price
is justified by the deliverable, not the other way around.

\subsection*{The Core Problem Connector Solves}

Every AI-forward company deploying LLM agents in 2025--2026 faces a version
of the same problem: \emph{the agent makes decisions and no one can fully
explain, constrain, or audit them at the infrastructure level.} Application-layer
logging exists but can be bypassed, edited, or simply insufficient for a
compliance auditor. The board, the compliance team, or the enterprise client
asks ``prove that your AI operated within policy'' and the engineering team's
answer is a manual log parse that takes days and still does not produce a
signed proof.

Connector installs a layer \emph{below} the application code that:

\begin{itemize}[leftmargin=1.5em, itemsep=3pt]
  \item Intercepts every agent memory read, write, and decision at the kernel
        level --- not at the app layer.
  \item Evaluates that action against an 11-layer policy chain before
        allowing or blocking it.
  \item Produces a cryptographically signed receipt of the outcome: what
        was decided, what data was accessed, which policy was applied, who
        authorised the capability, and when.
  \item Exposes that receipt in a format that a human auditor, a compliance
        tool, or a regulator can read without a developer present.
\end{itemize}

\noindent That is the product. Not ``AI governance infrastructure.'' Not
``an agent operating system.'' The product is: \textbf{a signed proof that
your AI agent did what it was supposed to do --- or didn't, and here is
exactly why it was stopped.}

\subsection*{Startup Pilot Tier: \$1,000--2,000\,/\,month $\times$ 6~months}

Target: post-seed AI startup or boutique AI consulting firm, 8--40 employees.

\begin{table}[H]
\centering
\renewcommand{\arraystretch}{1.55}
\begin{tabularx}{\textwidth}{@{}lX@{}}
\toprule
\textbf{Deliverable} & \textbf{What the customer has at the end of 6 months} \\
\midrule
Audit trail & Every agent decision in their primary workflow is logged at
  the kernel level with a CID-addressed, signed receipt. Not application
  logs. Not request/response pairs. A tamper-evident proof chain. \\
Policy enforcement & 3--5 custom policy rules authored to their specific
  workflow (e.g.\ ``agent cannot access PII fields without an explicit
  capability token'', ``agent budget cannot exceed N calls per session'').
  Enforced at the kernel level before the action executes. \\
Compliance starter & A one-page evidence export they can show to their
  Series~A investors, their own enterprise clients, or an early SOC2
  auditor: ``our AI agents operate under these policies, here is the
  proof from the last 30 days.'' \\
Deployment & A running Connector node in their infrastructure (Docker
  Compose single-node). The founders assist with setup on day one. \\
Integration & Their primary agent workflow (1--2 agent types) is integrated
  with the Connector kernel. SDK and worked examples provided. \\
Roadmap access & Direct Slack channel. Feature requests logged and
  prioritised. Response within business-hours on the same day. \\
\midrule
\textbf{Early pilot exclusive} & Price locked for 18 months post-pilot.
  Named as a design partner on the website (if agreed). Co-authorship of
  one integration pattern that ships in the official SDK. First access to
  the CLS authoring interface when it launches. \\
\bottomrule
\end{tabularx}
\caption{Startup pilot tier --- deliverables for \$1,000--2,000\,/\,month over 6~months}
\end{table}

\noindent The one-sentence pitch for this tier: \emph{``Your board asked
what happens when your AI makes the wrong decision. After six months with
Connector, your answer is a signed receipt showing exactly what it decided,
why, and which policy it was operating under --- provable without a developer
in the room.''}

\subsection*{Mid-Enterprise Pilot Tier: \$5,000--7,000\,/\,month $\times$ 6~months}

Target: Series~A or Series~B company, 40--200 employees, with 2--5 internal
AI agent workflows in production, compliance team already asking questions.

\begin{table}[H]
\centering
\renewcommand{\arraystretch}{1.55}
\begin{tabularx}{\textwidth}{@{}lX@{}}
\toprule
\textbf{Deliverable} & \textbf{What the customer has at the end of 6 months} \\
\midrule
Full multi-workflow coverage & All 2--5 of their production agent workflows
  instrumented. Not just one. Each workflow has its own policy set, its
  own receipt chain, and its own compliance export. \\
Compliance report export & Monthly one-click export: a PDF or CSV that
  their compliance team or external auditor accepts as evidence. Fields
  covered: agent identity, action taken, policy applied, data accessed,
  outcome (allowed / blocked), timestamp, cryptographic signature.
  Specifically structured for HIPAA or SOC2~CC6/CC7 evidence requirements. \\
Custom policy authoring & 8--12 policy rules authored together with the
  customer's compliance or engineering team. These are not template rules ---
  they are written for their specific data types, regulatory context, and
  agent behaviour patterns. \\
Real-time policy violation alerts & When an agent action is blocked by the
  kernel, an alert fires to their incident management system (PagerDuty,
  Slack, webhook). The alert includes: which agent, which rule, what it
  attempted, and what data it was trying to access. \\
Dedicated onboarding & Founders spend 2--3 days on-site or in intensive
  remote sessions to complete integration. Not a self-service process. \\
SLA for response & 4-hour business-hours response on any production issue
  via dedicated Slack channel. Founders are directly contactable. \\
Quarterly review & A 60-minute call each quarter reviewing agent behaviour
  trends, policy violation patterns, and recommended policy improvements. \\
\midrule
\textbf{Early pilot exclusive} & Annual contract option at 15\% discount
  (\$51K--\$71.4K/year vs.\ \$60K--\$84K at list price). Co-authored case
  study (published anonymously if required). Input rights on the enterprise
  roadmap for the following 12 months. First access to multi-tenant and
  multi-node features before general availability. \\
\bottomrule
\end{tabularx}
\caption{Mid-enterprise pilot tier --- deliverables for \$5,000--7,000\,/\,month over 6~months}
\end{table}

\noindent The one-sentence pitch for this tier: \emph{``You have five AI
agent workflows in production. Your compliance team asks quarterly: prove
they are operating within policy. After six months with Connector, that
question is answered with a one-click export --- every agent action, every
policy check, every data access, signed cryptographically --- structured
specifically for your HIPAA or SOC2 audit requirements. And when an agent
breaks a rule, the kernel stops it before the action happens and alerts your
on-call team within seconds.''}

\subsection*{What Early Pilots Get That Later Customers Never Will}

The first 10 pilot customers occupy a position that no subsequent customer
can buy:

\begin{enumerate}[leftmargin=1.8em, itemsep=4pt]
  \item \textbf{18-month price lock.} When Connector raises Seed~1 in Q4~2027
        and reprices to reflect the product's maturity, early pilots keep
        their pilot rate for 18 months from signing. Given the trajectory,
        this is likely a 30--50\% discount against future list price.

  \item \textbf{Direct founder access, permanently.} After Seed~1, the
        founders will be managing a team and a sales pipeline. Early pilot
        customers retain a Slack channel directly to the founders for the
        duration of their subscription. This is not offered to post-Seed
        customers.

  \item \textbf{Roadmap influence.} The first 10 pilots collectively shape
        30--40\% of the 2027 product roadmap. They are consulted before
        features are scoped, not after they are shipped.

  \item \textbf{Reference designation.} Customers who agree to be a named
        reference become part of the sales story for Seed~1. Their credibility
        raises capital at better terms, which means the company is more
        likely to survive and serve them well long-term. Being a reference
        is a direct investment in the vendor's success.

  \item \textbf{Co-authorship in the SDK.} The integration patterns the
        first 10 pilots use become the official SDK examples. Their workflow
        is the one every new customer sees first. That is structural product
        influence that persists beyond their contract.

  \item \textbf{Pre-announcement of enterprise features.} Before any new
        feature ships, early pilot customers are notified privately and
        given 2 weeks to test and provide feedback before general
        availability. They always have at least one version-cycle advantage
        over the general customer base.
\end{enumerate}

\begin{insightbox}{The pilot is not a discount. It is a co-investment.}
An early pilot customer is not getting a cheaper version of a finished
product. They are co-investing with their time, feedback, and money in
building a product that is shaped around their use case. The price reflects
that relationship. A customer who understands this signs. A customer who
thinks they are buying discounted enterprise software will be disappointed
by the rough edges and will churn. Qualify this distinction in every first
conversation.
\end{insightbox}

%%────────────────────────────────────────────────────────────
\section*{The Real Plan: Zero to \$30K MRR, Month by Month}
%%────────────────────────────────────────────────────────────

This section is not a projection. It is an operating plan for two specific
founders, one of whom has real financial pressure right now, in a codebase
that is architecturally strong but not yet customer-ready. Every week
matters. Every decision either extends runway or consumes it.

\begin{warnbox}{The honest starting condition --- March 2026}
One founder needs a salary and has real-life financial obligations that
cannot wait for a product-market fit story. This is not unusual and not
shameful. It is a constraint that must be treated the same way as a
technical constraint: acknowledged, planned around, and resolved as fast
as possible. The plan below is built around this constraint, not despite
it. The fastest path to a salary that covers real life is two paying pilots
by end of July. Everything else is in service of that goal.
\end{warnbox}

\subsection*{Before May: What Must Happen Right Now}

The window from now until the end of April is the only period where the
founders still have some breathing room. Use it for two things only:

\begin{enumerate}[leftmargin=1.8em, itemsep=4pt]
  \item \textbf{Generate one consulting income event.}\;
    One short-term project --- 2--3 weeks, \$3,000--5,000 --- extends
    personal runway by 6--10 weeks. One founder takes this. The other
    builds the demo. The goal is to arrive at May with \$8,000--12,000
    in combined pocket rather than \$0--5,000. This is not ideal but it
    is the difference between having a runway and not having one.
    Do not wait. Identify the project this week.

  \item \textbf{Complete the demo and deployment.}\;
    Docker Compose single-node. Works in one command. A script that walks
    through one governed agent workflow in under 10 minutes. Three worked
    examples in the integration doc. This is the entire Phase~0 deliverable.
    Nothing else gets built until this works reliably. Two founders, four
    weeks, full focus.
\end{enumerate}

\subsection*{Month-by-Month: May 2026 to November 2026}

\begin{table}[H]
\centering
\renewcommand{\arraystretch}{1.7}
\begin{tabularx}{\textwidth}{@{}llXll@{}}
\toprule
\textbf{Month} & \textbf{State} & \textbf{What happens} &
\textbf{MRR} & \textbf{Founder salary} \\
\midrule
May & \$0--5K pocket & Demo is ready. Outreach list of 75 CTOs built.
  Send first 40 messages. Target: 6--10 responses, 3--5 demos
  scheduled. One founder on consulting 1~day/week if needed.
  No revenue yet. & \$0 & \$0 from company \\
\midrule
June & Demos running & Run 8--12 demos. Qualify hard. Three serious
  conversations. Two LOI drafts in progress. Do not drop price.
  If a prospect says ``we like it but \$3,000 is too much''
  that is the wrong prospect. & \$0--\$3K
  (1--2 early invoices) & \$0--\$1.5K \\
\midrule
July & First LOIs & Two design partner agreements signed. First
  invoices paid within 7 days of signing. Pilot 1 and Pilot 2
  running. Concierge onboarding: founders do the setup. Weekly
  calls. Fix bugs same day. & \$6K--\$8K & \$3K--\$3.5K each \\
\midrule
August & 2 pilots live & Both design partners integrated and running.
  Focus: learn what they actually use, what breaks, what they
  would pay more for. Outreach continues: 40 more prospects.
  Target: 2 more closes before end of August. & \$10K--\$14K &
  \$4K--\$4.5K each \\
\midrule
September & 4--5 pilots & Four to five paying customers. Onboarding
  is faster because Phase~1 learnings are applied. One integration
  guide written and reused. Pilot feedback shapes top 3 product
  priorities. Angel warm-up conversations begin --- no ask yet,
  just the story. & \$14K--\$20K & \$5K each \\
\midrule
October & 6--7 pilots & Six to seven customers. One founder now spends
  60\%+ on customer integration issues --- this is the signal to
  start interviewing Tech~1. Pre-seed narrative taking shape:
  ``7 paying customers, \$18K--22K MRR, 90\% retention at
  3~months.'' & \$18K--\$24K & \$5K--\$5.5K each \\
\midrule
November & 7--10 pilots & Seven to ten paying customers. MRR
  \$25K--\$35K. Both founders on full living salary. Pre-seed
  conversations with 3--5 angels started. Tech~1 hire pipeline
  running. First customer ready to be a named reference. &
  \$25K--\$35K & \$6K each \\
\bottomrule
\end{tabularx}
\caption{Month-by-month plan --- May to November 2026, starting from near-zero cash}
\end{table}

\subsection*{The Personal Salary Resolution --- Month by Month}

This is the part that is usually not written down. It needs to be.

\begin{table}[H]
\centering
\renewcommand{\arraystretch}{1.6}
\begin{tabularx}{\textwidth}{@{}llXl@{}}
\toprule
\textbf{Month} & \textbf{MRR} & \textbf{What the founder can realistically take} &
\textbf{Gap to cover} \\
\midrule
March--April & \$0 & \$0 from company. Consulting income if arranged.
  Draw on savings. & \$3,500--5,000/mo personal \\
May--June & \$0--\$3K & \$0--\$1,500 from company. Still need
  savings or consulting. Tight. & \$2,000--5,000/mo \\
July & \$6K--\$8K & Company pays one founder \$3,000--\$3,500/month.
  The other either defers or uses savings for 4--6 more weeks. &
  \$0--\$2,000/mo per person \\
August & \$10K--\$14K & Both founders take \$4,000/month. Still lean
  but survivable in most cities. & \$0--\$1,000/mo \\
September & \$14K--\$20K & Both founders take \$5,000/month. Covers
  rent, food, basics. Not comfortable but stable. & Covered \\
November & \$25K--\$35K & Both founders take \$6,000/month. Real life
  bills are covered. Financial stress drops. & Covered with buffer \\
\midrule
\textbf{Q1 2027} & \$35K--\$50K & Post pre-seed close, both founders
  formalize \$6,500--7,500/month salary. First time it feels real. &
  Fully covered \\
\bottomrule
\end{tabularx}
\caption{Personal salary resolution timeline --- honest month-by-month}
\end{table}

\begin{warnbox}{July is the minimum viable month for personal survival}
If pilots 1 and 2 are not signed and invoiced by end of July, the
personal financial situation becomes a company crisis. One founder will
be forced to take outside income full-time. This is not failure --- it
is a signal that the ICP or the demo needs to change, not that the
company is dead. If July closes without revenue: spend two weeks
understanding why the demos did not convert, fix exactly one thing, and
restart outreach. Do not pivot the entire product. Fix the conversion, not
the product.
\end{warnbox}

\subsection*{What ``Design Partner'' Means in Practice}

Two design partners in July--August 2026 does not mean two customers who
use the platform occasionally. It means:

\begin{itemize}[leftmargin=1.5em, itemsep=3pt]
  \item A contract signed and a card charged before any integration work
        begins. Even \$1,000 upfront. Money before work.
  \item A weekly 30-minute call, on a fixed day, every week for the
        first 8 weeks. The founder drives the agenda. The customer shows
        up because they paid.
  \item A specific workflow already identified before the first call:
        ``your agent that does X --- that is what we are going to
        govern first.'' Not a general deployment. One workflow, done well.
  \item A clear success criterion agreed in writing at signing:
        ``at the end of 6 months, you will have a signed audit trail for
        every decision this agent makes, exportable to a one-page report
        for your compliance team.'' Not vague. Specific and measurable.
  \item The founders treat this customer's bugs as production incidents.
        Response within 2 hours. A fix within 24 hours or a workaround
        with an ETA. This is not normal startup product behaviour --- it
        is the reason they pay \$3,000--\$4,000/month for something that
        is not fully finished.
\end{itemize}

%%────────────────────────────────────────────────────────────
\section*{2027: Grounded Quarterly Plan}
%%────────────────────────────────────────────────────────────

Entering 2027, Connector has 7--10 paying pilots, \$25K--35K MRR, a
working product, and a real customer story. Q1~2027 is still a
single-track execution: close the pre-seed, bring Tech~1 in, renew the
2026 pilot cohort. From Q2~2027 the business runs on \textbf{three
commercial tracks simultaneously.} These tracks are not alternatives ---
they are parallel and permanent.

\subsection*{Three Commercial Tracks (Q2 2027 onwards --- all three run forever)}

The three tracks map directly to the licensing tier model already implemented
in code: \texttt{Tier::Growth}/\texttt{Business} for Track~1,
\texttt{Tier::Indie}/\texttt{Startup} for Track~2, and
\texttt{Tier::Enterprise}/\texttt{Core} (with \texttt{features\_override:
["hipaa", "soc2"]}) for Track~3.

\begin{table}[H]
\centering
\renewcommand{\arraystretch}{1.75}
\begin{tabularx}{\textwidth}{@{}clXXl@{}}
\toprule
\textbf{\#} & \textbf{Track} & \textbf{Exact customer type} &
\textbf{Why this track cannot be abandoned} & \textbf{Tier (code)} \\
\midrule
1 & \textbf{Survival Bridge}\newline Mid-startup \&\newline small business &
  Ecommerce AI support SaaS, AI agent startups, fintech tools --- the
  same companies from the 2026 ICP who renew and new companies of the same
  type. 20--50 employees. AI in production. Monthly pain. & This is the
  revenue floor. Even if every enterprise deal slips by 6 months, these
  customers keep the lights on and salaries paid. \textbf{Never deprioritise
  this track regardless of enterprise pipeline size.} &
  \texttt{Growth} \$500/mo\newline \texttt{Business} \$1,000/mo \\
\midrule
2 & \textbf{Self-Serve Small Biz}\newline General use &
  Small companies using AI internally: approvals, internal ops automation,
  small agent deployments. No sales touch. Signs up through the portal,
  activates a license, self-integrates. & Passive recurring revenue that
  compounds without founder time. 80 customers at \$200/month average =
  \$16K MRR with zero sales effort. Funds tech hires. &
  \texttt{Indie} \$150/mo\newline \texttt{Startup} \$250/mo \\
\midrule
3 & \textbf{Enterprise Regulated}\newline Long-term strategic &
  Healthcare (HIPAA pilot grant), fintech (SOC2/FCA), regulated ops
  teams. Uses \texttt{PilotGrant} with \texttt{features\_override:
  ["hipaa", "soc2"]} during onboarding. Slow sales cycle. High ACV. &
  These are the real long-term customers. 5 enterprise accounts at
  \$3K--\$4K/month = \$15K--\$20K MRR. NRR expands as they roll out
  more workflows. This is the Series~A story. &
  \texttt{Enterprise} \$3K/mo\newline \texttt{Core} \$4K/mo \\
\bottomrule
\end{tabularx}
\caption{Three permanent commercial tracks from Q2 2027 --- grounded in actual licensing tier code}
\end{table}

\begin{warnbox}{Track 1 is the survival constraint, not a stepping stone}
Connector's real long-term customers are enterprise in regulated segments.
But enterprise deals slip. A 3-month slip in one enterprise close is
normal. A 6-month slip happens regularly. If the company has abandoned
Track~1 to focus on enterprise, a slip means payroll stress and founder
salary cuts. Track~1 must always have active outreach and active closes
running in parallel with the enterprise motion, unconditionally. This
is not optional in good times and it is the only thing that matters in
bad times.
\end{warnbox}

\subsection*{Quarterly Plan with Three-Track Breakdown}

\begin{table}[H]
\centering
\renewcommand{\arraystretch}{1.7}
\begin{tabularx}{\textwidth}{@{}llXl@{}}
\toprule
\textbf{Quarter} & \textbf{MRR target} &
\textbf{Operational focus} & \textbf{Key event} \\
\midrule
Q1 2027 & \$35K--\$50K &
  \textit{Single track.} Close pre-seed SAFE (\$75K--\$125K). Tech~1 joins.
  Renew 2026 pilot cohort to annual contracts. Price for new Track~1
  customers rises to \$800--\$1,200/month. Track~2 self-serve portal goes
  live. No enterprise outreach yet --- build the compliance feature flag
  pipeline (\texttt{hipaa}, \texttt{soc2}) so Track~3 is ready. &
  Pre-seed closes \\
\midrule
Q2 2027 & \$50K--\$65K &
  \textit{Three tracks start.}
  Track~1: continue ecommerce/agent outreach, target 5 new Track~1 closes.
  Track~2: self-serve portal active, 30--50 small business sign-ups at
  Indie/Startup tier, no founder time needed per customer.
  Track~3: first 2--3 enterprise pilot conversations start (healthcare or
  fintech). Use \texttt{PilotGrant} with tier\_override = Enterprise and
  features = [hipaa/soc2]. Seed investor shortlist of 10--15 names built. &
  3-track begins;\newline first enterprise pilots \\
\midrule
Q3 2027 & \$65K--\$85K &
  Track~1: 25--30 active mid-startup customers at \$800--\$1K avg = \$20--\$30K.
  Track~2: 60--80 self-serve customers at \$200 avg = \$12--\$16K.
  Track~3: 3--4 enterprise pilots paying = \$9K--\$16K.
  Formal Seed~1 fundraise begins. Tech~1 carries product. Founders
  split 60\% / 40\% execution vs.\ fundraise. Track~1 must not stall during
  fundraise --- Tech~1 owns customer integration for Track~1 customers. &
  Seed~1 process begins \\
\midrule
Q4 2027 & \$75K--\$100K &
  Seed~1 closes (\$500K--\$1.5M at \$4M--\$6M post-money).
  Sales~1 hire approved (owns Track~3 enterprise pipeline).
  Both founders on \$7,500--\$9,000/month formal salary.
  Track split at year-end: T1 \$25--\$35K $+$ T2 \$15--\$20K $+$
  T3 \$20--\$30K = \$60K--\$85K MRR base entering 2028.
  Series~A story: NRR from Track~1 renewals $>$ 110\%,
  Track~3 ARR growing, Track~2 is passive proof of platform breadth. &
  Seed~1 closes \\
\bottomrule
\end{tabularx}
\caption{2027 quarterly plan with three-track commercial model from Q2 onwards}
\end{table}

\subsection*{What Makes 2027 Work or Fail}

The 2027 plan hinges on three things that cannot be bought or rushed:

\begin{enumerate}[leftmargin=1.8em, itemsep=5pt]
  \item \textbf{Annual renewal rate from 2026 pilots.}\;
    If fewer than 60\% of the November 2026 pilots renew into 2027,
    the Seed story is broken. The product is not retaining. The fix is
    not more sales --- it is diagnosing why the first cohort did not
    renew and fixing it before adding more customers. One churned pilot
    investigated honestly is worth more than five new ones chased blindly.

  \item \textbf{Tech~1 hire quality.}\;
    Tech~1 joins in Q1 2027. If this person cannot independently
    handle customer integration issues within 6 weeks of starting,
    the founders cannot give 40\% of their time to fundraising in Q3.
    The Seed process dies because product execution stalls. Hire slowly.
    One wrong Tech~1 sets the company back 4 months.

  \item \textbf{Founder bandwidth during Seed fundraise.}\;
    Raising Seed takes 60--90 hours of active founder time per month
    for 2--3 months. Product decisions slow down. Customer response
    times increase. This is unavoidable. The mitigation is:
    (a) give existing customers a heads-up that the founding team is
    in a fundraise period, (b) ensure Tech~1 has a documented
    escalation path for production issues, and (c) do not start the
    fundraise until the product is stable enough to run for 4 weeks
    without a founder intervention.
\end{enumerate}

\begin{insightbox}{The only metric that matters at the end of 2027}
At the end of Q4 2027, one number determines whether the company is on
track: the 12-month NRR of the 2026 pilot cohort. If it is above 110\%
--- meaning pilots renewed and some expanded their contract --- the
Series~A conversation in 2028 is credible. If it is below 90\%, no
amount of new logos fixes the story. Protect the 2026 cohort as if
they are the only customers you will ever have, because in 2027, they
effectively are.
\end{insightbox}

\subsection*{Honest Numbers: What Self-Sufficiency Actually Looks Like at Each Stage}

\begin{table}[H]
\centering
\renewcommand{\arraystretch}{1.55}
\begin{tabularx}{\textwidth}{@{}llXX@{}}
\toprule
\textbf{Period} & \textbf{MRR} & \textbf{Team burn} &
\textbf{Monthly net} \\
\midrule
Nov 2026 & \$30K & 2 founders at \$6K = \$12K total & \$+18K (reinvested
  in hiring pipeline) \\
Q1 2027  & \$42K & 2 founders + Tech~1 = \$20K & \$+22K \\
Q2 2027  & \$58K & + Tech~2 = \$26K & \$+32K \\
Q4 2027  & \$88K & + Sales~1, formalised salaries = \$40K &
  \$+48K (strong Seed~1 narrative) \\
Q4 2028  & \$140K & 8--10 people at \$80K--\$100K total & \$+40K--\$60K \\
Q4 2029  & \$160K--\$185K & 12--16 people at \$130K--\$150K &
  Cash-flow positive. No survival dependency on next round. \\
\bottomrule
\end{tabularx}
\caption{Team burn vs.\ MRR from November 2026 through end of 2029}
\end{table}

\begin{warnbox}{This plan requires that July 2026 actually closes revenue}
Every number in 2027 and 2028 is downstream of two pilots signing and
paying in July 2026. If that does not happen, the cascade shifts by
one quarter: August closes, October hits \$7K, January 2027 hits
\$30K. That is survivable. What is not survivable is not closing the
first pilot before October 2026. After that date, the founder's personal
financial situation forces a decision that may end the company. The
entire strategic plan lives or dies in a 30-day window in July. Treat
it accordingly.
\end{warnbox}
